\chapter{SQL Fundamentos}

\section[Sublanguages SQL]{\emj{🧠}~Sublanguages SQL}

SQL es un lenguaje declarativo: describes \emph{qué} quieres, no \emph{cómo} obtenerlo.
El motor decide el plan de ejecución óptimo.
Se divide en cinco sublanguages según el tipo de operación:

\begin{teoriabox}{Los cinco sublanguages}
\begin{itemize}
  \item \textbf{DDL} (Data Definition Language): estructura — \texttt{CREATE, ALTER, DROP, TRUNCATE}
  \item \textbf{DML} (Data Manipulation Language): datos — \texttt{INSERT, UPDATE, DELETE}
  \item \textbf{DQL} (Data Query Language): consultas — \texttt{SELECT}
  \item \textbf{DCL} (Data Control Language): permisos — \texttt{GRANT, REVOKE}
  \item \textbf{TCL} (Transaction Control): transacciones — \texttt{BEGIN, COMMIT, ROLLBACK, SAVEPOINT}
\end{itemize}
DDL es auto-commit en la mayoría de los DBMS: un \texttt{DROP TABLE} no puede hacerse \texttt{ROLLBACK} una vez ejecutado (excepto en PostgreSQL, donde DDL sí es transaccional).
\end{teoriabox}

\section[DDL --- definir estructura]{\emj{🧠}~DDL --- definir estructura}

El DDL define el esquema: qué tablas existen, qué columnas tienen y qué constraints las gobiernan.
Los constraints se pueden definir inline (junto a la columna) o como constraints de tabla separados.
\texttt{TRUNCATE} es más rápido que \texttt{DELETE} porque no genera undo log fila por fila, pero no puede ser filtrado con \texttt{WHERE}.

\begin{itemize}
  \item \texttt{CREATE TABLE} define el esquema; constraints inline o separados
  \item \texttt{ALTER TABLE} modifica columnas, agrega/quita constraints sin recrear la tabla
  \item \texttt{TRUNCATE} vacía la tabla rápido (no dispara triggers fila por fila)
  \item \texttt{DROP} elimina objeto permanentemente --- irreversible fuera de transacción
\end{itemize}

\section[DML + DQL --- manipular y consultar]{\emj{🧠}~DML + DQL --- manipular y consultar}

\texttt{SELECT} es el comando más complejo de SQL.
Su orden de escritura no es el orden de evaluación:
se escribe \texttt{SELECT} primero pero se evalúa casi al último.
Entender este orden es clave para saber qué puedes referenciar en cada cláusula.

\begin{tipbox}{Orden lógico de evaluación de SELECT}
\texttt{FROM} → \texttt{JOIN} → \texttt{WHERE} → \texttt{GROUP BY} → \texttt{HAVING} → \texttt{SELECT} → \texttt{ORDER BY} → \texttt{LIMIT}

Por eso no puedes usar un alias definido en \texttt{SELECT} dentro de \texttt{WHERE}: \texttt{WHERE} se evalúa antes de \texttt{SELECT}.
\end{tipbox}

\begin{lstlisting}[caption={DDL: CREATE, ALTER, DROP}]
-- Crear tabla con constraints
CREATE TABLE clientes (
    id        SERIAL PRIMARY KEY,
    email     VARCHAR(200) NOT NULL UNIQUE,
    nombre    VARCHAR(100) NOT NULL,
    activo    BOOLEAN DEFAULT TRUE,
    creado_en TIMESTAMP DEFAULT NOW()
);

-- Agregar columna
ALTER TABLE clientes ADD COLUMN telefono VARCHAR(20);

-- Agregar constraint después de crear
ALTER TABLE clientes ADD CONSTRAINT chk_email CHECK (email LIKE '%@%');

-- Eliminar columna
ALTER TABLE clientes DROP COLUMN telefono;
\end{lstlisting}

\begin{lstlisting}[caption={DML: INSERT, UPDATE, DELETE}]
-- INSERT con RETURNING (PostgreSQL)
INSERT INTO clientes (email, nombre)
VALUES ('ana@ejemplo.com', 'Ana López')
RETURNING id, creado_en;

-- UPDATE con condición
UPDATE clientes
SET    activo = FALSE
WHERE  creado_en < NOW() - INTERVAL '1 year'
  AND  activo = TRUE;

-- DELETE con subquery
DELETE FROM clientes
WHERE id IN (
    SELECT id FROM clientes
    WHERE activo = FALSE
    ORDER BY creado_en
    LIMIT 100
);
\end{lstlisting}

\begin{lstlisting}[caption={SELECT con JOINs}]
-- INNER JOIN: solo filas que coinciden en ambas tablas
SELECT e.nombre, d.nombre AS departamento
FROM   empleados e
INNER JOIN departamentos d ON e.id_depto = d.id;

-- LEFT JOIN: todas las filas de la izquierda, NULL si no hay match
SELECT e.nombre, d.nombre AS departamento
FROM   empleados e
LEFT JOIN departamentos d ON e.id_depto = d.id;

-- GROUP BY + HAVING
SELECT   d.nombre, COUNT(e.id) AS total, AVG(e.salario) AS promedio
FROM     empleados e
JOIN     departamentos d ON e.id_depto = d.id
GROUP BY d.nombre
HAVING   COUNT(e.id) > 3
ORDER BY promedio DESC;
\end{lstlisting}

\begin{entrevistabox}
\begin{enumerate}
  \item ¿Cuál es la diferencia entre \texttt{WHERE} y \texttt{HAVING}?
  \item ¿Qué diferencia hay entre \texttt{DELETE} y \texttt{TRUNCATE}?
  \item Explica los tipos de JOIN: INNER, LEFT, RIGHT, FULL OUTER.
\end{enumerate}
\end{entrevistabox}
